\phantomsection\addcontentsline{toc}{section}{Tranformers para clasificación de textos}
\section*{Tranformers para clasificación de textos}

Ya hemos hablado acerca de que la arquitectura Transformer representó una ruptura fundamental con
los modelos recurrentes y convolucionales tradicionales dentro del Deep Learning. Eliminar la 
dependencia explícita de la secuencialidad y cambiarla por un mecanismo llamado ``auto-atención'' que
permite procesar todos los tokens en paralelo, vovlió a esta arquitectura en toda una revolución dentro
del NLP, y el DL en general. 

En lugar de propagar los estados ocultos de manera recurrente como en las RNN o los LSTM/GRU, el 
Transformer calcula relaciones directas entre todas las posiciones de una secuencia mediante pesos 
de atención que cuantifican la relevancia contextual de cada token respecto a los demás. Matemáticamente, 
estos trabajan con tres proyecciones que ya mencionamos en la sección pasada de Transformers, pero
para el caso generativo, las matrices de consulta $Q$, clave $K$ y valor $V$:

\[
    Q = XW_Q. \;\;\;\;\;\;\; K = XW_K, \;\;\;\;\;\;\; V = XW_V
\]

donde $W_Q, W_K, W_V \in \mathbb{R}^{d \times d}$, son matrices de parámetros. La atención para cada
token $i$ se calcula como una combinación ponderada de los valores $V$, donde los pesos derivan de
una similitud escalar entre las consultas y las claves:

\[
    \text{Attention}(Q,K,V) = \text{Softmax} \Big(\frac{QK^T}{\sqrt{d_k}}\Big) V
\]

Este cálculo define la contribución de cada token al contexto del token actual en la secuencia. Al 
combinar varias cabezas de atención, el modelo puede capturar simultáneamente distintos tipos de relaciones
tanto semánticas como sintácticas. 

\subsection{Clasificación de texto con Transformers}

En tareas de clasificación, como lo es análisis de sentimineto, categorización temática o detección
de polaridad, el Transformer se usa como codificador contextual. Cada token produce un embedding contextualizado
$h_i$, y un vector de resumen de la secuencia se pasa a una capa densa para estimar la probabilidad 
de la clase:

\[
    \hat{y} = \text{Softmax} \big(W_c h_{\text{[CLS]} + b_c}\big)
\]

La función de pérdida por excelencia es la entropía cruzada entre la distribución predicha $\hat{y}$ 
y la verdadera $y$:

\[
    \mathbb{L}_{CE} = - \sum_{k=1}^{K} y_k \text{log} \hat{y}_k
\]

Durante el fine-tuning, las capas del transformer se pueden actualizar de manera parcial o total, acorde
al tamaño del corpus y los recursos disponibles. En la práctica, modelos como BETO, BERT-base o DistilBERT 
logran estados del arte en clasificación textual, esto gracias a aus capacidad de capturar relaciones semánticas
a largo plazo sin recurrencia explícita. 

\subsection{Contraste con el enfoque generativo}

La distinción entre el uso generativo y discriminativo de los Transformers radica, principalmente,
en el objetivo del enternamiento y el flujo de la información. Podemos ver esto en la siguiente tabla:

\begin{table}[h!]
    \centering
    \caption{Comparación entre modelos de lenguaje generativos y discriminativos.}
    \label{tab:model_comparison}
        \begin{tabular}{lll}
        \toprule
        \textbf{Aspecto} & \textbf{Generativo (GPT, LLaMA)}        & \textbf{Discriminativo (BERT, BETO)}             \\ \midrule
        Tipo de atención & Causal (unidireccional)                 & Bidireccional                                    \\
        Objetivo         & Predicción del siguiente token          & Predicción de tokens    \\
        Salida           & Texto continuo                          & Etiqueta de clase                                \\
        Métrica típica   & Perplejidad (PPL)                       & Accuracy, F1-score                               \\
        Uso típico       & Generación de texto, diálogo            & Clasificación, análisis de sentimientos          \\ \bottomrule
        \end{tabular}
\end{table}

En la práctica, ambos modelos comparten la misma base estructural, bloques de atención, normalización
y proyecciones lineales, pero difieren en cómo restringen el acceso al contexto y en el tipo de 
tarea supervisada que se les plantea. Mientras que los modelos generativos como GPT o LLaMA se entrenan 
para predecir la próxima palabra en la secuencia, los modelos tipo BERT/BETO internalizan el contexto 
completo para producir embeddings discriminativos que pueden ser reutilizados para tareas de fine-tuning 
eficinete. Este cambio de paradigma logró unificar, en una misma arquitectura, el aprendizaje de representaciones 
linguisticas, siendo así el estado del arte actual y con lo que habrían soñado los fundadores de esta
ciencia. 

